\chapter*{The field behind the symmetry}
\addcontentsline{toc}{chapter}{The field behind the symmetry}

Every finite group can act faithfully on a finite set and a rational vector space. The inverse Galois problem requires a further realization: a finite Galois field extension of the rational numbers whose full automorphism group is the prescribed group. Its field multiplication and arithmetic base are part of the question.

The recovered Operation Komorebi manuscript organizes fixed objects through subgroup averages, expresses normality and compatibility in a regular module, and transports traces, norms and descent data into a resolution language. These structures describe an existing extension and clarify what an inverse construction must produce. A faithful coordinate action alone does not supply the field.

The reader begins with quadratic, biquadratic and cubic fields. They make automorphisms explicit and distinguish a root field from its splitting field. A quartic combines exact modular factors with the external arithmetic cycle theorem. Independent variables then realize any finite group over its invariant field, leaving the arithmetic base change as a separate obligation.

The later chapters develop subgroup projectors, regular specialization and embedding problems. A counterexample shows how a faithful group representation can lose subgroup information retained by the regular group algebra action. A weak lift is distinguished from a surjective lift. The small examples can be checked directly. Universal implications have written proofs or named external inputs.

\begin{claimboundary}
This edition proves the specified elementary field and projector statements, explains conditional realization routes and records bounded exact calculations. It does not establish inverse Galois realization over $\Q$ for every finite group. No proof assistant certification is claimed.
\end{claimboundary}

\section*{How to read the evidence}

A written argument supports a universal statement under declared hypotheses. A primary source supplies an external theorem through an exact source contract. A finite calculation checks its specified objects and range. A model illustrates a distinction. An open obligation identifies missing input. These evidence roles remain separate.

The bundle preserves four recovered inputs byte for byte: the expanded Galois manuscript package, its separate embedding ledger, and the Frobenius packet findings and laboratory package. That packet multiplication work concerns finite rings and transforms, not inverse Galois existence. This reader does not claim to have independently audited every theorem in the larger historical manuscript.

The literature is dated 9 October 2026. Classical existence results are source attributed. Recent $M_{23}$ work is recorded as a versioned external preprint whose computations are outside this audit. The monograph, visual guide and source bundle share the same claim boundary and identify the next proof obligations.

